\section{Data Audit Study}\label{sec:audit}
To validate the existence of evaluation bias in public IoT datasets (captured by the model in Section~\ref{sec:background}), we audited two datasets that were released for IoT anomaly detection~\cite{gummadi2024xaiiot}. The audit uses the released files without modification, and the dataset statistics are presented in Table~\ref{tab:audit}.

\subsection{MEMS Vibration Dataset}
The MEMS dataset was collected on a motor testbed with a balancing disk~\cite{gummadi2024xaiiot,abdallah2023sensors}. MEMS accelerometers record the acceleration along three axes at a sampling rate of 10\,Hz. The motor operates at ten rotational speeds between 100 and 600 revolutions per minute, and a mass mounted on the disk creates mechanical imbalance and thereby three health conditions, namely normal operation, near-failure, and failure. The released file contains 21,576 samples with three features and three classes, and it contains no session or speed column. The file encodes the conditions as the integer labels 1, 2, and 3, which contain 10,965, 5,773, and 4,838 samples. Following the class designation of~\cite{gummadi2024xaiiot}, we read these labels as normal, near-failure, and failure. MEMS accelerometers cost about \$8, compared with about \$1305 for piezoelectric sensors, but they offer lower sampling rates~\cite{gummadi2024xaiiot}.

We therefore cut the file into contiguous blocks of identical labels. This yields 15 blocks with a median length of 1,423 rows, which we use as an inferred grouping for block protocols. The z-axis shows a lag-1 autocorrelation of 0.38, whereas the x- and y-axes show 0.06 and 0.12, so that neighboring samples are partly redundant.

\subsection{N-BaIoT Traffic Dataset}
N-BaIoT contains traffic of nine commercial IoT devices, among them doorbells, a thermostat, a baby monitor, and security cameras, that were infected with the Mirai and Gafgyt botnets~\cite{meidan2018nbaiot,antonakakis2017}. The repository that accompanies~\cite{gummadi2024xaiiot} provides a subset of 20 features for each device. Each device has either six or eleven classes of equal nominal size. We audit all nine devices and use device 9, which has eleven classes and 209,000 rows, for the experiments.

Table~\ref{tab:audit} lists the audit results, in which cleaning removes exact duplicates and the rows of ambiguous vectors, so that every N-BaIoT device loses 41.8 to 48.3\% of its rows. The smallest class keeps only 9 to 23 unique rows, and two or three classes fall below 50 unique rows. MEMS, in contrast, is almost free of duplicates (0.1\%) and retains 4,832 unique rows in its smallest class.
\begin{table*}[t]
\centering
\caption{Audit of the released datasets. Bold: best N-BaIoT value per column. Smallest class: unique rows of the smallest class after cleaning.}
\label{tab:audit}
\scriptsize
\setlength{\tabcolsep}{3pt}
\begin{tabular}{@{}lrrrrrrrr@{}}
\toprule
Dataset & Classes & Rows & Unique rows & Removed (\%) & Smallest class & Ambiguous vectors & Rows in them (\%) & Accuracy ceiling \\
\midrule
mems & 3 & 21,576 & 21,551 & 0.1 & 4,832 & 6 & 0.1 & 1.000 \\
\midrule
nbaiot:1 & 11 & 220,000 & 116,635 & 47.0 & 17 & \textbf{5} & 43.4 & 0.689 \\
nbaiot:2 & 11 & 132,000 & 76,786 & \textbf{41.8} & 11 & 6 & 40.3 & 0.719 \\
nbaiot:3 & 6 & 120,000 & 68,937 & 42.6 & \textbf{23} & \textbf{5} & \textbf{39.1} & \textbf{0.833} \\
nbaiot:4 & 11 & 220,000 & 123,680 & 43.8 & 12 & 8 & 43.3 & 0.689 \\
nbaiot:5 & 11 & 220,000 & 113,767 & 48.3 & 22 & 6 & 47.4 & 0.648 \\
nbaiot:6 & 11 & 220,000 & 114,335 & 48.0 & 14 & 7 & 47.0 & 0.652 \\
nbaiot:7 & 6 & 120,000 & 66,652 & 44.5 & 19 & 7 & 39.3 & \textbf{0.833} \\
nbaiot:8 & 11 & 220,000 & 127,915 & 41.9 & 17 & 10 & 39.9 & 0.724 \\
nbaiot:9 & 11 & 209,000 & 117,073 & 44.0 & 9 & 11 & 40.5 & 0.718 \\
\bottomrule
\end{tabular}
\end{table*}


The ambiguity is more severe than the duplication of records in these datasets. On each device, 5 to 11 feature vectors occur under more than one label, and they cover 39.1 to 47.4\% of all rows. On device 9, 11 vectors cover 40.5\% of the rows. These vectors account for 99.5, 100, and 99.9\% of the rows of classes 1, 3, and 4, respectively, and a single vector occurs under eight labels (Table~\ref{tab:ambclass} in Appendix~\ref{app:audit}). Equation~\eqref{eq:ceiling} then bounds the accuracy at 0.718 on device 9 and between 0.648 and 0.833 across all devices. For MEMS, only 0.1\% of the rows are ambiguous, and the bound equals 1.000.

\textbf{Generalizability of the study:} The audit covers one sensor dataset and all nine devices of the traffic dataset. The duplication and ambiguity statistics are similar across devices, which suggests that they are a property of the released 20-feature subset and not of device 9. Random forests and ensembles reach accuracies between 0.64 and 0.83 on the same files in~\cite{gummadi2024xaiiot}, which coincides with the bounds above. We therefore hypothesize that the lower accuracy of the 20-feature subset, in comparison with all 115 features, reflects information that the subset discards and not a limited quality of the models. We could not test this hypothesis, because the full feature set is not part of the released files.
